% Solutions of 21-divide-and-conquer.
\section*{Chapter \ref{ch:b1:divide-and-conquer} --- Divide and Conquer}

\begin{solution}{exo:b1:divide-and-conquer:1}
The 3 is larger than both 1s and the 2; the 4 than the second 1 and the 2; the 5 than the 2; the 9 than the 2 and the 6:
$3 + 2 + 1 + 2 = 8$. The left half $[3, 1, 4, 1]$ has 3 inversions, the right half $[5, 9, 2, 6]$ has 3, and the
top-level merge counts the 2 pairs across the middle, $(3, 2)$ and $(4, 2)$.
\end{solution}

\begin{solution}{exo:b1:divide-and-conquer:2}
$a_1 = 12$, $a_0 = 34$, $b_1 = 56$, $b_0 = 78$. $z_2 = 12 \times 56 = 672$, $z_0 = 34 \times 78 = 2652$, $z_1 = |12 -
34| \cdot |56 - 78| = 22 \times 22 = 484$, and $s = (-1)(-1) = 1$. The middle term is $672 + 2652 - 484 = 2840$ (indeed
$12 \times 78 + 34 \times 56 = 936 + 1904$), and $672 \cdot 10^4 + 2840 \cdot 10^2 + 2652 = 7\,006\,652$.
\end{solution}

\begin{solution}{exo:b1:divide-and-conquer:3}
Pivot 5: below $[2, 4, 1]$, above $[7, 9, 8]$; rank 5 is above, at rank $5 - 4 = 1$ there. Pivot 8: below $[7]$, one key
equal; rank 1 is the pivot: the answer is 8, the key of rank 5 of $1, 2, 4, 5, 7, 8, 9$.
\end{solution}

\begin{solution}{exo:b1:divide-and-conquer:4}
\code{covers} is false for the widths 0 to 9 (at width 9 the bands $[0, 9]$, $[10, 19]$ and $[20, 29]$ are needed) and
true from 10 on ($[0, 10]$ and $[12, 22]$). The dichotomy on $[0, 20]$ tests 10 (true), 5, 8 and 9 (false), and returns
10 after four tests.
\end{solution}

\begin{solution}{exo:b1:divide-and-conquer:5}
The halves give $\delta \approx 4.123$, and the strip around $x = 5$ holds $(4, 6)$, $(5, 8)$, $(7, 2)$ and $(9, 7)$. In
the $y$ order they are $(7, 2)$, $(4, 6)$, $(9, 7)$, $(5, 8)$: the closest pair, $(4, 6)$ and $(5, 8)$, at $\sqrt5
\approx 2.236$, is separated by $(9, 7)$, which is far from both in $x$. Neighbours in the $y$ order are not enough: the
proposition allows up to seven points between, and the loop must go on while the $y$ distance is below $\delta$.
\end{solution}

\begin{solution}{exo:b1:divide-and-conquer:6}
Near $10^5$, two consecutive doubles differ by $2^{16 - 52} = 2^{-36} \approx 1.5 \times 10^{-11}$, more than the
tolerance. When $lo$ and $hi$ are consecutive doubles, their mean rounds to one of them, the loop assigns it, and
nothing changes: the output shows $mid = hi$. Repairs: stop when $mid$ equals $lo$ or $hi$; or use a relative tolerance,
$hi - lo > 10^{-12} \cdot hi$; or run a fixed number of halvings (about 110 here cover the 50 bits of the initial width
and the 53 bits of precision).
\end{solution}

\begin{solution}{exo:b1:divide-and-conquer:7}
\omcode{code/bachelor-1/21-divide-and-conquer/exercises/max_subarray.py}{4}{19}{Largest slice sum}
A best slice lies in the left half, in the right half, or across the middle, where it is a best slice ending at $mid -
1$ followed by a best slice starting at $mid$; both are found by one scan. So $T(n) = 2T(n/2) + \Theta(n) = \Theta(n
\log n)$. The test compares with Kadane's linear scan and with all slices on 200 random lists.
\end{solution}

\begin{solution}{exo:b1:divide-and-conquer:8}
\omcode{code/bachelor-1/21-divide-and-conquer/exercises/majority.py}{4}{13}{Majority key}
If a key $c$ held at most half of each half, it would hold at most $\lfloor mid/2 \rfloor + \lfloor (n - mid)/2 \rfloor
\leq n/2$ positions: a majority key of the list is a majority key of one half. So it is one of the at most two
candidates returned by the halves, and counting each costs $\Theta(n)$: $T(n) = 2T(n/2) + \Theta(n) = \Theta(n \log n)$.
\end{solution}

\begin{solution}{exo:b1:divide-and-conquer:9}
For numbers of $d$ digits, \code{multiply} takes the smallest power $2^k \geq d$, so $2^k < 2d$ and $3^k = (2^k)^{\log_2
3} < (2d)^{\log_2 3} = 3 d^{\log_2 3}$. With a cutoff $2^j$, $K(2^k) = 3K(2^{k-1})$ for $k > j$ and $K(2^j) = 4^j$, the
school method's count: $K(2^k) = 3^{k-j} 4^j$.
\end{solution}

\begin{solution}{exo:b1:divide-and-conquer:10}
With groups of three, at least half of the $n/3$ medians are at most the pivot, and each of their groups holds two keys
at most the pivot: about $n/3$ keys on each side, so the recursive call receives up to $2n/3$, and $T(n) \leq T(n/3) +
T(2n/3) + an$. Now $1/3 + 2/3 = 1$: no level of the \omterm{def:b1:asymptotic-analysis:recurrence}{recursion tree} is smaller than the one above. Take $T(n) = T(n/3) +
T(2n/3) + an$ and suppose $T(j) \geq a j \log_3 j$ for $j < n$. Then
\[ T(n) \geq a\tfrac n3 \log_3 \tfrac n3 + a\tfrac{2n}{3} \log_3 \tfrac{2n}{3} + an = an \log_3 n - a\tfrac n3 -
a\tfrac{2n}3 + a\tfrac{2n}3 \log_3 2 + an \geq an \log_3 n, \]
so $T(n) = \Omega(n \log n)$; with groups of three the method is not proved linear.
\end{solution}

\begin{solution}{exo:b1:divide-and-conquer:11}
\omcodes{code/bachelor-1/21-divide-and-conquer/exercises/skyline.py}{5-30}{The skyline}
Merging two outlines walks both lists once, keeping the current height of each, and records the larger one wherever it
changes: $\Theta(n)$ for outlines of $O(n)$ points. Hence $T(n) = 2T(n/2) + \Theta(n) = \Theta(n \log n)$. The test
compares the outline's height at every integer abscissa with the highest building standing there.
\end{solution}

\begin{solution}{exo:b1:divide-and-conquer:12}
\omcode{code/bachelor-1/21-divide-and-conquer/exercises/four_sum.py}{7}{21}{Four keys with a given sum}
The table keeps, for each sum of two keys, the pair $(k, m)$ with that sum and the largest $k$. If a solution $i < j < k
< m$ exists, then when the loops reach $(i, j)$, the sum $a[k] + a[m]$ is in the table with a pair whose first position
is at least $k > j$: a solution is returned. Conversely a returned quadruple has the target sum and increasing positions.
The table takes $\Theta(n^2)$ steps, its sorting $O(n^2 \log n)$, and the $n(n - 1)/2$ binary searches $O(n^2 \log
n)$. The test compares with all quadruples on random lists.
\end{solution}

\begin{solution}{pb:b1:divide-and-conquer:1}
\textbf{1.} $n^2$, one for each pair of digits: 1\,048\,576 for $n = 1024$.

\textbf{2.} Cell $t$ receives the products $a_i b_j$ with $i + j = t$: at most one $j$ for each $i$, so at most $n$
products, each at most $9 \times 9 = 81$.

\textbf{3.} One pass over $2n$ cells, with a carry at most $(81n + \mathit{carry})/10 < 9n$: $\Theta(n)$ steps on
numbers of a few bytes. Each step replaces $r_t + \mathit{carry}$ by its last digit and carries the rest to cell $t + 1$,
which keeps $\sum_t r_t 10^t$ unchanged; at the end every cell is a digit.

\textbf{4.} $a, b < 10^n$, so $ab < 10^{2n}$.

\textbf{5.} From 9412 to 39\,270 to 167\,668: by about 4.2 per doubling; the model predicts 4.

\textbf{6.} A million, $1000^2$; the program pads to 1024 digits and makes 1\,048\,576.

\textbf{7.} $ab = a_1 b_1 10^{2m} + (a_1 b_0 + a_0 b_1) 10^m + a_0 b_0$.

\textbf{8.} $(a_1 - a_0)(b_1 - b_0) = a_1 b_1 - a_1 b_0 - a_0 b_1 + a_0 b_0$; move the two middle terms to the other side.

\textbf{9.} $0 \leq a_0, a_1 < 10^m$, so $|a_1 - a_0| < 10^m$.

\textbf{10.} The digits of $|x - y|$, as a list of the same length, and the sign of $x - y$. Between two numbers of the
same length, the first digit that differs from the top decides which is larger.

\textbf{11.} The loop makes $\sum_t r_t 10^t$ equal to the product. \code{divmod} writes any integer, negative or not, as
$10q + d$ with $0 \leq d \leq 9$, so each step keeps the sum and leaves a digit; the product is in $[0, 10^{2n})$, so
the last carry is 0 and the cells are its digits.

\textbf{12.} $K(1) = 1$ and $K(n) = 3K(n/2)$.

\textbf{13.} $K(2^0) = 1 = 3^0$, and $K(2^{k+1}) = 3K(2^k) = 3 \cdot 3^k = 3^{k+1}$.

\textbf{14.} $3^k = 2^{k \log_2 3} = (2^k)^{\log_2 3} = n^{\log_2 3}$.

\textbf{15.} $3^{10} = 59\,049$ against 1\,048\,576: about 18 times fewer.

\textbf{16.} The slices, the two differences, the combine loop and \code{normalise} each take $\Theta(n)$ steps. So
$W(n) = 3W(n/2) + \Theta(n)$, and since $3 > 2^1$ the master theorem gives $\Theta(n^{\log_2 3})$: the leaves dominate,
and the additions do not change the order.

\textbf{17.} $3^{k-j} 4^j$; for $n = 1024$ and $c = 32$, $3^5 \cdot 4^5 = 248\,832$.

\textbf{18.} The cutoff 1: each level handed to the school method multiplies the count by $4/3$. A larger cutoff is
used because a call, with its lists and its carries, costs far more than a digit product: the time is not the count.

\textbf{19.} 12.3, 10.8, 8.1, 6.9, 5.1, 2.8, 2.0, 1.4 for $n = 16$ to 2048. The factor falls by about 0.7 per doubling
at the end, close to the $3/4$ expected when one time grows like $3^k$ and the other like $4^k$.

\textbf{20.} With factors of about $1.4 \times 0.75 \approx 1.06$ at 4096 digits and 0.8 at 8192: from about 5000
digits.

\textbf{21.} With cutoff 8, from 256 digits (1534 against 1773 microseconds); with cutoff 32, from 64 digits (99 against
113).

\textbf{22.} $29\,815 / 9772 \approx 3.05$ and $9772 / 3204 \approx 3.05$: close to 3, where the school method grows by
4.

\textbf{23.} $70 \times 30 = 2100$ bits, and $2100 \log_{10} 2 \approx 632$ decimal digits.

\textbf{24.} \emph{Karatsuba's count}: Karatsuba's method multiplies two numbers of $2^k$ digits with exactly $3^k$
digit products, that is $n^{\log_2 3} \approx n^{1.585}$ for $n = 2^k$ (questions 12 to 14), the counts printed by the
program in \cref{out:b1:divide-and-conquer:karatsuba}; its whole cost is $\Theta(n^{\log_2 3})$ (question 16).
\end{solution}
